% !PRINTED 558
\weberpage{558}
%<<CONT

\gesperrt{II. Every number in $R(r)$ which admits the substitution $(r, r_e)$, and hence also every number of the field $R(\eta)$, can be represented as a homogeneous linear function of the periods $\eta, \eta_1 \ldots \eta_{e-1}$.}

For by I. we can put every number of $R(r)$ in one way into the form

\weq{13}{\varphi(r) = \sum a_h r_h}

where the coefficients $a_h$ are rational numbers, and $h$ runs through a complete system of residues with respect to the modulus $n - 1$. But we have
\[
\varphi(r_e) = \sum a_h r_{h+e} = \sum a_{h-e} r_h,
\]
and therefore, if $\varphi(r) = \varphi(r_e)$ is to hold, we must have $a_{h-e} = a_h$ for every $h$; hence also $a_h = a_{h+e} = a_{h+2e} \ldots$ Accordingly (13) can be put into the form
\[
\varphi(r) = \sum_{0,\,e-1}^{h} a_h r_h + \sum_{0,\,e-1}^{h} a_h r_{h+e} + \cdots + \sum_{0,\,e-1}^{h} a_h r_{h+(f-1)e},
\]
or

\weq{14}{\varphi(r) = a\eta + a_1\eta_1 + \cdots + a_{e-1}\eta_{e-1},}

which was to be proved. It is further to be remarked here that, if $\varphi$ has integer coefficients, the coefficients of the $\eta_h$ also become integers.

Since every number of the field $R(\eta)$ can be represented in one and only one way in the linear form (14), we call the system of quantities $\eta, \eta_1, \eta_2 \ldots \eta_{e-1}$ a \gesperrt{basis of the field $R(\eta)$}.

If one knows the representation of the product of each two periods $\eta_h \eta_k$ in the form (14), one can from it, by repeated application, represent all rational functions of the periods $\eta$ in the same form. Likewise one can also easily form the equation whose roots are the $e$ quantities $\eta$, if one applies the procedure to the product

\weq{15}{F_e(x) = (x - \eta)(x - \eta_1) \cdots (x - \eta_{e-1})}

This equation can also be represented in determinant form.

Namely, let us put

\weq{16}{\eta\eta_h = a_{0,h}\eta + a_{1,h}\eta_1 + \cdots + a_{e-1,h}\eta_{e-1},}

in which the $a_{0,h}, a_{1,h} \ldots a_{e-1,h}$ are integers, about whose calculation further details will follow later. If one sets up this equation for $h = 0, 1, \ldots e - 1$ in the form:
%CONT>>
